% Generated by scripts/assemble.py; corrections belong in source/corrections.json.
\examhead{71}{大连理工大学 2026 年数学分析试卷}{71}
% p072-b03; source PDF page 72
\par\Needspace{5\baselineskip}
\blockanchor{p072-b03}
\majorquestion{一}简答题.\par
% p072-b04; source PDF page 72
\blockanchor{p072-b04}
\question{1}试构造仅在 $x = 0$ 和 $x = 2026$ 两点可微的函数 $f(x)$, 而在其余个点皆不连续.\par
% p072-b05; source PDF page 72
\blockanchor{p072-b05}
\question{2}(可能有误) 计算极限 $\lim_{x\to 0^{+}}\frac{\sqrt{e^x - 1} - e\sqrt[x]{x}}{\sqrt[x]{x}}.$\par
% p072-b06; source PDF page 72
\blockanchor{p072-b06}
\question{3}设 $f(x)$ 可微, 且有 $\lim_{x\to +\infty}\frac{xf(x) + f'(x)}{x} = L\in \mathbb{R}$, 证明: $\lim_{x\to +\infty}f(x) = L$.\par
% p072-b07; source PDF page 72
\blockanchor{p072-b07}
\question{4}设  $\{a_{n}\}$  是正数列, 证明:  $\lim_{n\to\infty}\sqrt[n]{a_{1}^{n}+a_{2}^{n}+\cdots+a_{n}^{n}}=\sup\{a_{n}\}$ .\par
% p072-b08; source PDF page 72
\blockanchor{p072-b08}
\question{5}证明 $\sin (x^2 + y^2)$ 在 $\mathbb{R}^2$ 上不一致连续.\par
% p072-b09; source PDF page 72
\blockanchor{p072-b09}
\question{6}计算极限 $\lim_{\substack{x \to \infty \\ y \to \infty}} \frac{2x - 3y}{x^2 - 2xy + 3y^2}$.\par
% p072-b10; source PDF page 72
\blockanchor{p072-b10}
\question{7}计算 $f(x) = \int_{0}^{x} \frac{\mathrm{d}t}{\sqrt{1 - t}}$ 的幂级数展开式.\par
% p072-b11; source PDF page 72
\blockanchor{p072-b11}
\question{8}计算极限 $\lim_{x\to \infty}\int_0^\pi \max \{1,y\} \sin^2 (xy)\mathrm{d}y.$\par
% p072-b12; source PDF page 72
\blockanchor{p072-b12}
\question{9}证明: $\sin x > \frac{2}{\pi} x, 0 < x < \frac{\pi}{2}$.\par
% p072-b13; source PDF page 72
\blockanchor{p072-b13}
\question{10}$\{a_{n}\}$ 收敛当且仅当对任给的正整数 $p$, 有 $\lim_{n\to \infty}|a_{n + p} - a_n| = 0$, 此命题是否正确? 为什么?\par
% p072-b14; source PDF page 72
\par\Needspace{5\baselineskip}
\blockanchor{p072-b14}
\majorquestion{二}计算题.\par
% p072-b15; source PDF page 72
\par\Needspace{6\baselineskip}
\blockanchor{p072-b15}
\question{1}计算曲线积分\par
\nopagebreak[4]
% p072-b16; source PDF page 72
\blockanchor{p072-b16}
\displayformula{\int_ {L} \left(y ^ {2} + z ^ {2}\right) \mathrm{d} x + \left(z ^ {2} + x ^ {2}\right) \mathrm{d} y + \left(x ^ {2} + y ^ {2}\right) \mathrm{d} z.}
% p072-b17; source PDF page 72
\blockanchor{p072-b17}
其中曲线 $L$ 为球面 $x^{2} + y^{2} + z^{2} = 1 (x \geq 0, y \geq 0, z \geq 0)$ 与三坐标平面的交线, 它的定向使得球面的上侧在曲线左侧.\par
% p072-b18; source PDF page 72
\blockanchor{p072-b18}
\question{2}设 $D \subset \mathbb{R}^2$ 由 $x^2 + y^2 = 2y$ 和 $y \geq 1$ 围出, 计算二重积分 $\iint_{D} \frac{2y - y^2 - x^2}{y} \, \mathrm{d}x \, \mathrm{d}y$.\par
% p072-b19; source PDF page 72
\blockanchor{p072-b19}
\question{3}设 $a > 0, b > 0$, 计算积分 $\int_0^{+\infty} \frac{\ln(a^2 + x^2)}{b^2 + x^2} \mathrm{d}x$.\par
% p072-b20; source PDF page 72
\par\Needspace{5\baselineskip}
\blockanchor{p072-b20}
\majorquestion{三}证明题.\par
% p072-b21; source PDF page 72
\blockanchor{p072-b21}
\question{1}设 $f(x)$ 在 $[a, b]$ 上连续，且 $f(a) < 0 < f(b)$. 证明：存在 $c \in (a, b)$，使得 $f(c) = 0$，且 $f(x) > 0$，$x \in (c, b]$ 时.\par
% p072-b22; source PDF page 72
\par\Needspace{6\baselineskip}
\blockanchor{p072-b22}
\question{2}设 $D \subset \mathbb{R}^2$ 为区域, $u(x, y), v(x, y)$ 在 $D$ 上连续可微, 且在 $D$ 上满足\par
\nopagebreak[4]
% p072-b23; source PDF page 72
\blockanchor{p072-b23}
\displayformula{\frac {\partial u}{\partial x} = \frac {\partial v}{\partial y}, \frac {\partial u}{\partial y} = - \frac {\partial v}{\partial x}, u ^ {2} (x, y) + v ^ {2} (x, y) = 2026.}
% p072-b24; source PDF page 72
\blockanchor{p072-b24}
证明： $u(x,y),v(x,y)$ 在 $D$ 上是常数.\par
% p072-b25; source PDF page 72
\blockanchor{p072-b25}
\question{3}定义 $f(x) = \int_0^x \cos \frac{1}{t} \, \mathrm{d}t, x \in \mathbb{R}$，证明：$f(x)$ 在 $x = 0$ 可导，且 $f'(0) = 0$.\par
% p073-b00; source PDF page 73
\par\Needspace{6\baselineskip}
\blockanchor{p073-b00}
\question{4}设 $f(x)$ 是 $\mathbb{R}$ 上以 $2\pi$ 为周期的连续函数, 定义\par
\nopagebreak[4]
% p073-b01; source PDF page 73
\blockanchor{p073-b01}
\displayformula{f _ {n} (x) = \frac {1}{a _ {n}} \int_ {- \pi} ^ {\pi} f (t) \left(\cos \frac {t - x}{2}\right) ^ {2 n} \mathrm{d} t, x \in \mathbb {R}, n = 1, 2, 3, \dots .}
% p073-b02; source PDF page 73
\blockanchor{p073-b02}
其中  $a_{n}=\int_{-\pi}^{\pi}\cos^{2n}\left(\frac{x}{2}\right)\mathrm{d}x,\ n=1,2,\cdots$  ，证明：函数列  $\{f_{n}(x)\}$  在 R 上一致收敛于  $f(x)$ .\par
% p073-b03; source PDF page 73
\blockanchor{p073-b03}
\question{5}设 $f(x)$ 在 $[1, +\infty)$ 上非负递减，且 $\int_{1}^{+\infty} x^{a} f(x) \, \mathrm{d}x$ 收敛，证明：$\lim_{x \to +\infty} x^{a+1} f(x) = 0$.\par
